\section{A fixed testing family and its affine equations}\label{app:functions}
The preceding identities give an evaluated finite certificate cone.
This section specifies its testing functions, including both state signs,
and proves that the cone is closed. Its admission statement is precisely
Theorem~\ref{thm:cap}; no complementary testing region is included here.

First take $u,v\geq0$ and use the common state sign to make $w\geq0$.
There are two closed charts, $\sigma w\zeta\geq0$, with $\sigma=+1,-1$.
All functions below are divided by $\sqrt{D_v}$. Their squared arguments
are the following fixed polynomials:
\[
 q_0=1,\quad q_1=uv,\quad q_2=uw\zeta,\quad q_3=vw\zeta
       \quad(\sigma=+1),\qquad
 q_0=1,\quad q_1=uv,\quad q_4=-uvw\zeta
       \quad(\sigma=-1).
\]
These weights are nonnegative on their respective charts. All source
radicals in a row are coefficients; for example
$\sqrt{UV}=x^{-1/4}y^{-1/4}\sqrt{uv}$ and
$\sqrt{Uw\Zc}=x^{-1/4}q_z^{-1/4}\sqrt{uw\zeta}$.
Thus no source parameter occurs in a testing-function argument.

\subsection*{A finite support prescription}
For a row $f$ with just one testing radical $\sqrt{q_k}$, let
$\operatorname{supp}_k f$ be the monomials in $u,v,t,r,w,\zeta$ that
remain after removing that radical and expanding. In taking this
support, source-only coefficients are formal symbols: possible
specializations to zero never delete a monomial. This convention makes
the family independent of the source, including at $n=0$.

Here is a complete finite prescription for the rows whose supports
are used. The ten noncircuit rows of $f$ in the main text are included,
as are the three circuit factors on the positive chart or the five
factors
\[
 Uw,\quad Vw,\quad\Zc,\quad UV\Zc,\quad2\sqrt{-UVw\Zc}
\]
on the negative chart. The latter identity is
$\Ps=U^2w^2+V^2w^2+\Zc^2+U^2V^2\Zc^2-4UVw\Zc$.
The retained squares in Appendix~\ref{app:parent} have row supports
\begin{gather*}
 N,\ uN,\ vN,\ v(N+p\bc),\ uM,\ vM,\ uvM,\quad
 V\Qh,\ V\Bh,\ V(w-UV\Zc),\ V\Qh-\Rh,\ \Qh,\ \Rh .
\end{gather*}
Every coefficient of these squares is nonnegative on the source domain.

Let $\mathbf A,\mathbf B,\mathbf E$ be the displayed square-root
vectors of $\Aop,\Bop,\E$ in \eqref{eq:elementary}--\eqref{eq:opposite};
let $\mathbf B_0=(\Rh,U\Qh,U(w-UV\Zc))$.
Let $\mathbf W_\sigma,\mathbf W^\pi_\sigma$ be the vectors in
\eqref{eq:W19} and \eqref{eq:negativeW}; on the positive chart their
constant positive blocks are $G_7,1,G_3$. Include every row of
\begin{gather*}
 \mathbf W_\sigma,\quad U\mathbf A,\quad
 U\mathbf W^\pi_\sigma,\quad (U^2-1)(U\Zc-\Bh),\quad
 (\alpha/B)R_u-\alpha uR,\quad\mathbf B_0,\\
 v\mathbf A,\quad u^2v\mathbf A,\quad
 v\mathbf B,\quad u^2v\mathbf B,\quad
 \sqrt T\,v(\alpha+c_0u^2)R,\quad
 \mathbf A,\quad\mathbf B,\quad v\mathbf E,\quad u\mathbf E,
 \qquad c_0=BTCH.
\end{gather*}
This list contains every positive remainder of the radial, $J_\alpha$,
$O$, signed-mass and helper identities. In particular, $\bar a$ is
split into its $v^2$ and $u^4v^2$ terms before taking roots; no
$\sqrt{\alpha^2+c_0^2u^4}$ is introduced.

The eight event coordinates are represented by the following rows,
with the squared rows in each group summed:
\[
\begin{array}{c|l}
 \lambda&W+vX\\
 n&X,\ vX\\
 m&vw-t,\ u(vw-t)\\
 C&r,\ ur,\ vr,\ uvr\\
 A&Y,\ vY\\
 D&r+v\zeta,\ u(r+v\zeta)\\
 \alpha&vY-Z\\
 \beta&t,\ ut,\ vt,\ uvt
\end{array}
\]
Here $W,X,Y,Z$ are the testing polynomials in \eqref{eq:rotation}.
These rows allow arbitrary nonnegative event additions literally.

For the four residual pivots take the rows of $\Ps,\Qh^2,\bc^2,\Rh^2$,
respectively. Their nonconstant loss supports are
\[
\begin{array}{c|l}
 \Ps&(0,2),(0,4),(2,0),(2,2),(4,0),(4,2),(6,2)\\
 Q&(0,2),(0,4),(2,0),(4,0),(4,4)\\
 E&(0,4),(2,0),(2,2),(4,0),(6,2)\\
 R&(0,2),(2,0),(4,2).
\end{array}\tag{\(*\)}
\]
For each pivot row $f$ and each $(a,b)$ in its row of $(*)$, include
$u^{a/2}v^{b/2}f$ in the base supports. Also form a separate block with
rows $\sqrt{1-P_{ab}}\,f$, where
$P_{ab}=(2^{12}u)^a(2^{25}v)^b$.
These rows implement both coefficient signs in the residual identity
in the main text. Constant zero coefficients may be retained.

Finally the two signs of the three matrix gates have respective
polynomial supports
\[
 \begin{array}{c|c|l}
 (a,b)&\hbox{weight}&\hbox{monomial vector}\\
 (1,1)&1\pm P_{11}&(t,r,vw,u\zeta,v\zeta,uw)\\
 (0,2),(2,2)&1\pm P_{ab}&(r,v\zeta,uw).
 \end{array}
\]
All rows participating in each nonconstant core coefficient lie in
the indicated span; the circuit rows do not participate.

For each chart and each radical class, collect the union of its base
supports into one monomial vector. For each scalar gate collect the
supports belonging to that exponent and class into a separate vector.
The matrix gates are the six vectors just displayed. A block with
vector $m_b$ and weight $q_b$ represents
\[
        D_v^{-1}q_b\,m_b^TQ_bm_b,\qquad Q_b\succeq0.
\]
This support prescription is an explicit finite construction, using
only polynomial expansion and collection of distinct monomials.
Its evaluated vectors and rational weights are in
\texttt{fixed\_family.json}; the full row list is in
\texttt{root\_catalog.json}. They give
\[
\begin{array}{c|rrr|r}
 \hbox{chart}&\hbox{base}&\hbox{scalar gates}&\hbox{matrix gates}&\hbox{sum}\\
 +&63&56&24&143\\
 -&58&49&24&131.
\end{array}
\]
There are $49$ Gram blocks, $2269$ upper-triangular Gram variables,
and total positive-semidefinite order $143+131+9=283$, including the
nine scalar source-sign blocks. A constant positive Gram block such
as $G_7$ is transferred to its monomial vector by a congruence with
source-only coefficients. Sums of such contributions remain positive.
Thus grouping rows does not change any identity or require a mixed inverse.

\subsection*{The rational affine system and its scope}
Let $T_c=D_vS_c$ after \eqref{eq:rotation}. Its nine coefficient columns are
\begin{gather*}
 T_\lambda=(W+vX)^2,\quad T_n=D_vX^2,\quad
 T_m=D_u(vw-t)^2,\quad T_C=D_uD_vr^2,\quad T_A=D_vY^2,\\
 T_D=D_u(r+v\zeta)^2,\quad T_\alpha=(vY-Z)^2,\quad
 T_\beta=D_uD_vt^2,\qquad
 T_d=(T_D+T_\alpha+T_\beta-D_vZ^2)/2.
\end{gather*}
For each chart impose the polynomial coefficient equations
\[
       \sum_{b\ {
m in\ chart}}q_bm_b^TQ_bm_b
          =\sum_{i=1}^9 c_iT_i .
\]
The file \texttt{affine\_system.json} records these $554$ equations
($277$ per chart): $2707$ Gram coefficients and $132$ source coefficients.
Indices start at zero; off-diagonal entries are actual symmetric Gram
entries and hence contribute twice. Every record is rational.

These equations and positive-semidefinite constraints define a cone
$K_{283}$ in the nine source coefficients. The displayed identities
and coefficient reserves prove that it contains every native source
in \eqref{eq:domain}, including $n=0$, and the eight event rays.
Conversely, each of its sources has $S_c\geq0$ for all states on the
closed testing cap: all block weights are nonnegative there.
To cover signed angles, apply
\[
 (u,v,t,r,w,\zeta)\longmapsto
 (a u,b v,t,ab r,bw,a\zeta),\qquad a,b\in\{-1,1\}.
\]
All nine $T_i$ are invariant. Reducing the angles to nonnegative
values and then fixing the common state sign reaches one of the two
charts. The testing axes and $w\zeta=0$ require no division and are
included by the same formulas. This is a certificate cone for the
stated cap, rather than a claim of adequacy on the remaining tests.

\subsection*{A fixed positive grid proves closedness}
Every monomial in $m_b$ has individual degree at most
$(4,3,1,1,1,1)$ in $(u,v,t,r,w,\zeta)$.
On either chart choose the $320$ rational points in
\[
 \left\{\frac{i u_0}{6}:1\leq i\leq5\right\}
 \times\left\{\frac{k v_0}{5}:1\leq k\leq4\right\}
 \times\{1/3,2/3\}^{3}\times\sigma\{1/3,2/3\}.
\]
All weights $q_b$ are strictly positive at these points. Distinct
monomials with the displayed individual degree bounds are independent
on this Cartesian grid by successive univariate interpolation.
Consequently the rational matrix
\[
       H_b=\sum_{a\ \rm in\ grid}
              \frac{q_b(a)}{1+v(a)^2}\,m_b(a)m_b(a)^T
\]
is positive definite for every block.

Suppose $c^{(k)}\in K_{283}$ converges to $c$. Evaluate its identity
on these finitely many points and sum. Each summand from each Gram
block is nonnegative, so $\operatorname{tr}(H_bQ_b^{(k)})$ is bounded
above by the sum of the target values. These target sums converge
because they depend linearly on $c^{(k)}$. Since $H_b\succ0$, every
$Q_b^{(k)}$ is bounded. There are finitely many blocks, so a common
subsequence converges to positive-semidefinite matrices satisfying
the same coefficient equations at $c$. Hence $K_{283}$ is closed.
This argument also handles dependencies between different blocks;
it needs only independence within each explicitly chosen monomial vector.
